\section{Casos de aplicación y evolución cronológica}

\subsection{Contexto de los benchmarks y línea de tiempo}

\begin{knowledgebox}{Línea de tiempo de la evolución de Q-Learning}
\begin{itemize}
    \item 1992: Watkins demuestra la convergencia teórica de Q-Learning.
    \item 2013: Mnih et al. logran control end-to-end en Atari con DQN.
    \item 2015-2017: Double DQN, Dueling, Prioritized Replay, Rainbow y otras variantes van perfeccionando el método.
    \item Después de 2016: DQN se extiende a tareas de control continuo como MuJoCo y Roboschool, y empiezan a popularizarse el RL multi-agent y el entrenamiento adversarial.
\end{itemize}
\end{knowledgebox}

Al proponer un benchmark siempre hay que elegir entre usar un entorno sencillo para obtener reproducibilidad o un entorno complejo para aproximarse a las aplicaciones reales. El repaso del curso de Stanford subrayó el trade-off entre el control que ofrece el primero y la representatividad del segundo.

\subsection{Lecciones de los casos}

\begin{importantbox}{Revisión del caso de Atari}
En la tarea Atari Breakout, el Agent extrae información de movimiento mediante frame stacking, mientras que frame-skip mejora la eficiencia del entrenamiento. Al modificar el reward shaping (por ejemplo, al aumentar la recompensa por golpear un brick), la estrategia del agent cambió de forma notable, lo que muestra que ajustar las reward scales puede provocar fácilmente estrategias nuevas.
\end{importantbox}

\begin{warningbox}{Sesgo de observación en aplicaciones reales}
En las tareas de robótica, el ruido de los sensor reales y las diferencias con los sensor simulados pueden hacer que falle la transferencia sim-to-real. Aunque se alcance una puntuación alta en simulación, es indispensable aplicar domain randomization y calibrar con datos reales antes del despliegue.
\end{warningbox}

\subsection{Resumen de la sección}

Esta sección combina la evolución de los benchmarks con casos clásicos para recordarnos que, aun cuando dominemos los detalles del algoritmo, debemos cuidarnos de la brecha de información entre la simulación y el mundo real.
